\section*{Capítulo \ref{ch:g8:colors-spectra} --- Los colores y los espectros de la luz}

\begin{solution}{exo:g8:colors-spectra:1}
La banda continua de colores escondida en la \omterm{def:g1:light-and-shadows:source}{luz} blanca, ya
desempaquetada; la \omterm{prop:g8:colors-spectra:mixture}{dispersión} es el desempaquetado --- cada color
desviado su propia cantidad. El rojo se desvía lo mínimo; el violeta,
lo máximo.
\end{solution}

\begin{solution}{exo:g8:colors-spectra:2}
El arcoíris se forma en la mitad del cielo opuesta al Sol --- la \omterm{def:g1:light-and-shadows:source}{luz}
del Sol entra en las gotas de lluvia por detrás del cazador. Cada gota
hace de prisma (con un rebote de \omterm{def:g5:mirrors-reflection:reflection}{espejo} dentro).
\end{solution}

\begin{solution}{exo:g8:colors-spectra:3}
Desempaquetar: el prisma abre la \omterm{def:g1:light-and-shadows:source}{luz} blanca en abanico formando el
\omterm{prop:g8:colors-spectra:mixture}{espectro}. Empaquetar: el disco de arcoíris que gira se
desdibuja hasta parecer blanquecino --- y el segundo prisma pliega un
\omterm{prop:g8:colors-spectra:mixture}{espectro} de vuelta en un solo \omterm{def:g6:light-propagation:ray}{haz} blanco.
\end{solution}

\begin{solution}{exo:g8:colors-spectra:4}
La hierba absorbe casi toda la mezcla y devuelve sobre todo verde. La
página devuelve casi todo --- blanca; la pizarra no devuelve casi nada
y lo absorbe todo --- negra.
\end{solution}

\begin{solution}{exo:g8:colors-spectra:5}
\omterm{def:g1:light-and-shadows:source}{Luz} blanca: la bufanda devuelve su amarillo --- bufanda amarilla.
\omterm{def:g1:light-and-shadows:source}{Luz} azul: la oferta (azul) se encuentra con la respuesta
(amarillo) en un cruce vacío --- la bufanda se ve negra.
\end{solution}

\begin{solution}{exo:g8:colors-spectra:6}
La manzana solo puede devolver rojo, y el foco azul no trae ninguno:
negra. El pañuelo devuelve lo que le llegue --- si le ofrecen azul,
devuelve azul: un pañuelo azul.
\end{solution}

\begin{solution}{exo:g8:colors-spectra:7}
Rojo, verde y azul. Rojo $+$ verde: amarillo. Verde $+$ azul: cian.
Azul $+$ rojo: magenta. Los tres: blanco.
\end{solution}

\begin{solution}{exo:g8:colors-spectra:8}
Solo tiras diminutas rojas, verdes y azules, todas encendidas --- no
hay blanco por ninguna parte. La mezcla ocurre en tu ojo, demasiado
lejos para separar las tiras.
\end{solution}

\begin{solution}{exo:g8:colors-spectra:9}
Las luces suman lo que traen, y el juego completo de colores
juntos es blanco --- la suma completa la mezcla. Las pinturas absorben,
cada una restando su parte de la \omterm{def:g1:light-and-shadows:source}{luz}, y el juego completo de
restas no deja casi nada --- oscuridad por comité.
\end{solution}

\begin{solution}{exo:g8:colors-spectra:10}
El \omterm{def:g2:air-around-us:air}{aire} roba preferentemente el azul; todo el cielo recibe ese azul
robado y lo reenvía hacia abajo como la cúpula azul del día; el resto
rojo anaranjado sobrevive a la carretera rasante del
\omterm{def:g1:day-and-night:daynight}{atardecer}. A la \omterm{def:g2:moon-in-the-sky:moon}{Luna} eclipsada solo la ilumina
\omterm{def:g1:light-and-shadows:source}{luz} que ha rozado el borde de la Tierra --- filtrada por el
\omterm{def:g1:day-and-night:daynight}{atardecer} en todo el contorno del \omterm{def:g4:solar-system:planet}{planeta} ---, así que brilla con
ese mismo superviviente rojo.
\end{solution}

\begin{solution}{exo:g8:colors-spectra:11}
Las lámparas cálidas rojizas ofrecían poco verde; la respuesta
verdadera de la chaqueta --- sobre todo verde con algo de rojo, que da
una mezcla apagada --- solo podía enseñar su parte amiga del rojo, y
eso se leyó como un marrón cálido. La oferta completa de la \omterm{def:g1:light-and-shadows:source}{luz} del día
dejó hablar a la respuesta verde: engaño de una \omterm{def:g1:light-and-shadows:source}{luz} tacaña,
condenado por la regla del cruce.
\end{solution}

\begin{solution}{exo:g8:colors-spectra:12}
Letras azules, fondo negro. Con las luces blancas de sala, las letras
azules devuelven su azul mientras el fondo negro no devuelve nada:
letras de un azul brillante sobre negro --- legibles. Bajo el baño rojo
profundo, a las letras azules no se les ofrece nada que puedan devolver
y el fondo negro no devuelve nada, como siempre: negro uniforme,
pancarta desaparecida. Unas letras blancas traicionarían el final: el
devolvedor universal responde incluso a un baño rojo con letras de un
rojo brillante, legibles contra el fondo negro.
\end{solution}

\begin{solution}{pb:g8:colors-spectra:1}
\textbf{1.} Blanca --- los tres primarios a plena potencia suman la
mezcla completa.
\textbf{2.} Amarilla.
\textbf{3.} Camisa blanca: devuelve la oferta --- amarilla. Banda azul:
no le ofrecen azul --- negra. Vestido amarillo: la respuesta amarilla
coincide con la oferta amarilla --- amarillo brillante.
\textbf{4.} Fondo azul. La tela blanca devuelve lo que le mande la
diseñadora: un solo fondo, todos los colores a la carta --- el paralelo
de la pantalla blanca del cine.
\textbf{5.} Un abrigo verde: bajo el baño rojo, su respuesta verde no
se encuentra con ningún verde --- villano negro; con \omterm{def:g1:light-and-shadows:source}{luz} blanca,
la oferta completa deja hablar al verde --- verde intenso.
\textbf{6.} Una túnica blanca --- el devolvedor universal: roja bajo el
foco rojo, verde bajo el verde, azul bajo el azul y blanca bajo los
tres; siempre lo más brillante del escenario, siempre ``de aspecto
blanco'' en relación con la escena.
\textbf{7.} Cualquier iluminación \emph{sin} rojo: el foco azul, el
foco verde o los dos juntos --- si no le ofrecen rojo, la bufanda no
devuelve nada y se funde con la cortina negra. En el instante en
que se les une el foco rojo, la oferta contiene exactamente lo que la
bufanda responde: arde en rojo, y el truco muere.
\textbf{8.} El dorado, tratado como amarillo cálido, responde con
generosidad al rojo y al verde, y apenas al azul. El preajuste cargado
de azul le ofrecía poco que pudiera devolver: letras turbias y
apagadas. La mezcla completa del mediodía dejó arder el amarillo.
\textbf{9.} Franja roja: rojo brillante. Naranja y amarilla: versiones
rojizas y apagadas de sí mismas (sus respuestas incluyen algo de rojo).
Verde: negra. Azul y violeta: negras (la violeta, apenas de un oscuro
de brasa). El arco se derrumba en una ruina roja y negra.
\textbf{10.} Reviven las franjas roja y azul; la violeta brilla
débilmente (responde al azul y a algo de rojo); la naranja y la
amarilla solo enseñan sus flojas partes rojas; la franja verde sigue
negra --- la \omterm{def:g1:light-and-shadows:source}{luz} magenta no contiene nada de verde.
\textbf{11.} Solo la \omterm{def:g1:light-and-shadows:source}{luz} blanca --- la mezcla completa --- enseña
las seis franjas: cada franja devuelve su propio color, y solo una
oferta que contenga los seis puede ser respondida por los seis. Un foco
verde, por potente que sea, enciende una franja y asesina cinco.
\textbf{12.} Por ejemplo: ``Las luces mezcladas \emph{suman}: el
rojo, el verde y el azul construyen todos los colores y, juntos, el
blanco. Y todos los objetos \emph{restan}: devuelven solo su parte de
lo que trajo la \omterm{def:g1:light-and-shadows:source}{luz} --- no le ofrezcas nada que pueda devolver y
se vestirá de negro''.
\end{solution}
